\chapter{Error Control}

\section{Checksum}
% Checksum is a simple error-detection method where data is divided into fixed-size segments, added using binary addition with wrap-around, and the 1’s complement of the sum is sent as the checksum. The receiver verifies integrity by checking if the total (including checksum) equals all 1s.

Checksum is an error-detection technique used to ensure data integrity during transmission. It helps detect errors by appending a checksum value derived from the data. The receiver verifies integrity by checking if the total (including checksum) equals all 1s.

\vspace{1em}
\textbf{Advantages:}
\begin{itemize}[parsep=0.25em]
   \item Simple to compute and implement in hardware/software.
   \item Efficient for detecting common transmission errors.
   \item Low computational overhead, and widely used in protocols (TCP/IP, UDP)
\end{itemize}

\section{Practice Problems}
\textbf{\textit{Question:}}
Consider, the data unit to be transferred is 10010011 10010011 10011000 01001101. Checksum is an efficient error detection method to detect whether there is an error. Suppose, 8 bit checksum is used in this case. Now use this method on the data unit given above and answer the following:
Briefly explain the above-mentioned method and why it is a good choice.
Calculate the checksum and the exact data that will be transmitted.

\textbf{Solution:}

\begin{minipage}[t]{0.45\textwidth}
In sender side,\\[0.5em]
\begin{tabular}{rl@{}l}
      & 100100&11\\[-0.5em]
      & 100100&11\\[-0.5em]
      & 100110&00\\[-0.5em]
  $+$ & 010011&01\\\hline
   10 & 000010&11\\[-0.5em]
      &       &10\\\hline
      & 000011&01
\end{tabular}

\vspace{1em}
$\therefore$ Checksum is: 11110010 \ \ (1's compliment)\\[0.5em]
$\therefore$ The data will be transmitted as,\\
10010011 10010011 10011000 01001101 \rgbB{11110010}
\end{minipage}%
\hfill\vrule\hfill%
\begin{minipage}[t]{0.45\textwidth}
In receiver side,\\[0.5em]
\begin{tabular}{rl@{}l}
      & 100100&11\\[-0.5em]
      & 100100&11\\[-0.5em]
      & 100110&00\\[-0.5em]
      & 010011&01\\[-0.5em]
  $+$ & 111100&10\\\hline
   10 & 111111&01\\[-0.5em]
      &       &10\\\hline
      & 111111&11
\end{tabular}

\vspace{1em}
$\therefore$ Since the sum is 11111111, which is all 1s,\\
the data is correct.
\end{minipage}
